\documentclass[10pt,a4paper]{amsart}
\usepackage[margin=19mm]{geometry}
\usepackage{amsmath,amssymb,mathtools}
\usepackage[scr=rsfs,cal=euler]{mathalfa}
\usepackage{xfrac}
\usepackage[unicode,hidelinks,bookmarks=false]{hyperref}
\newcommand{\ie}{i.e.\ }
\newcommand{\cf}{cf.\ }
\newcommand{\resp}{respectively\ }
\newcommand{\Subsection}{\subsection}
\providecommand{\nobreakdash}{\nobreak\hbox{-}}
\setlength{\emergencystretch}{2em}
\multlinegap0pt
\usepackage{bm}
\usepackage{bbm}
\usepackage{dsfont}
\usepackage{xspace}
\usepackage{cancel}
\usepackage{mathtools}
\usepackage{cleveref}
\usepackage{amsthm}
\makeatletter
\@ifundefined{theorem}{\newtheorem{theorem}{Theorem}}{}
\@ifundefined{lemma}{\newtheorem{lemma}[theorem]{Lemma}}{}
\makeatother
\newcommand{\R}{\mathbb R}
\newcommand{\ang}[1]{\langle#1\rangle}
\newcommand{\cE}{\mathcal E}
\newcommand{\ceil}[1]{\lceil#1\rceil}
\newcommand{\var}{\textup{Var}}
\renewcommand{\P}{\mathbb{P}}
\renewcommand{\ge}{\geqslant}
\renewcommand{\le}{\leqslant}
\title[Improved lower bound for hypercube edge slicing]{Improved lower bound for hypercube edge slicing}
\author{Lisa Sauermann, Zixuan Xu}
\date{Source: arXiv:2510.16592v1 (CC BY 4.0)}
\begin{document}
\maketitle

\noindent\emph{Source and licence.} Improved lower bound for hypercube edge slicing. Version arXiv:2510.16592v1 (CC BY 4.0), distributed under the Creative Commons Attribution 4.0 International licence, as recorded in the arXiv licence metadata for this exact version and shown on the abstract page https://arxiv.org/abs/2510.16592v1. Full rights notice: licenses/arxiv-2510.16592-CC-BY-4.0.txt.\par\smallskip

\noindent\textit{Editorial scope.} Counted unit: the Lemma whose source label is \texttt{lem:almost-indep} in the article (source0.tex lines 784--787, source label \texttt{lem:almost-indep}), delivered together with the article's own complete original proof of it (source0.tex lines 789--829, one \texttt{proof} environment of 41 lines). No mathematics is written, repaired or re-derived: statement and proof are the article's own text line for line. Required context retained uncounted: lem:continuous-littlewood-offord (lines 290--293). Every definition, display and label of those lines is retained unchanged. Omitted from this package and not redistributed: 1--289, 294--783, 788--788, 830--1080. Nothing inside a retained excerpt is omitted or altered: every retained body is delivered exactly as the article's lines are written, so no omission receipt of any kind accompanies this package. Citations: the retained excerpts carry no citation command at all, so no bibliography entry of the article is transcribed and no reference is redistributed. Reference adaptation: none was needed and none was made; every reference inside the delivered unit resolves inside it, so no unresolved reference or citation is produced. Printed slips kept literal: no printed word, symbol, hypothesis or number of the counted statement or of its proof was altered. Layout and class: the article's own class is \texttt{amsart} with packages including amsmath, amssymb, amsthm, amsfonts, tikz-cd, mathrsfs, mathtools, stmaryrd, enumitem, algorithmic, float, comment, tikz, hyperref; the wrapper is the lane's pinned \texttt{amsart} editorial wrapper at 10pt on A4, which restates the theorem-like environments the retained text needs and adds the article's own macro definitions that the retained text needs (\texttt{\textbackslash P}, \texttt{\textbackslash R}, \texttt{\textbackslash ang}, \texttt{\textbackslash cE}, \texttt{\textbackslash ceil}, \texttt{\textbackslash ge}, \texttt{\textbackslash le}, \texttt{\textbackslash var}), with extra wrapper packages (bm, bbm, dsfont, xspace, cancel, mathtools, cleveref). The delivered numbering is the unit's own. Assets: no retained span contains a figure environment, a graphics-inclusion command, a PDF-page inclusion, a table or a TikZ picture; no image file is copied into this package, the package declares no asset dependency and no image file is redistributed with it. Licence: version arXiv:2510.16592v1 of this article is distributed under the Creative Commons Attribution 4.0 International licence, as recorded in the arXiv licence metadata for this exact version and shown on the abstract page https://arxiv.org/abs/2510.16592v1. A full rights notice is delivered at licenses/arxiv-2510.16592-CC-BY-4.0.txt. Characters: every retained excerpt is pure ASCII, so the delivered engine, XeLaTeX with its default Latin Modern text font, reports no missing character.
\begin{lemma}\label{lem:continuous-littlewood-offord}
    Let $a_1,\dots, a_n, b_1,\dots, b_n$ be real numbers with $a_i<b_i$ for all $i\in [n]$, and let $X_i\sim U[a_i, b_i]$ be independent random variables for all $i\in [n]$. Define $X = \sum_{i = 1}^n X_i$. Then for any $t>0$ and $b\in \R$, we have
    \[\P[|X - b| < t]\le \frac{2t}{\sqrt{\var(X)}}.\]
\end{lemma}

\begin{lemma}\label{lem:almost-indep}
    For distinct $i,j\in [\ell]$ with $|\ang{v_i, v_j}|\le 0.9$ and $t\in \{1,2,4,\dots,2^{\ceil{\log_2 m}}\}$, we have 
    \[\P[i\text{ bad}\text{ and }\cE_1(j,t)]\le  \frac{40000t}{\rho_0^2\sqrt{S_j}}\log m.\]
\end{lemma}

\begin{proof}
We first show that for any given outcomes of $\alpha_h$ for all $h\in [\ell]\setminus\{i,j\}$, we have
\begin{equation}\label{eq:conditional-on-alpha-h}
\P[i\text{ bad}\text{ and }\cE_1(j,t) \mid (\alpha_h)_{h\in [\ell]\setminus\{i,j\}}]\le  \frac{2000t}{\rho_0^2}\log m.
\end{equation}

To show \eqref{eq:conditional-on-alpha-h}, let us fix arbitrary outcomes of $\alpha_h$ for all $h\in [\ell]\setminus\{i,j\}$. Let $A = \ang{v_i, v_j}$. By assumption we have $|A|\le 0.9$. Note that 
\[\ang{X_0, v_i} = \rho_0\Big(\alpha_i + \alpha_j\cdot A + \sum_{h \in [\ell]\setminus\{i,j\}}\alpha_h\ang{v_h, v_i}\Big)\]
and 
\[\ang{X_0, v_j} = \rho_0\Big(\alpha_j + \alpha_i\cdot A + \sum_{h \in [\ell]\setminus\{i,j\}}\alpha_h\ang{v_h, v_j}\Big),\]
and that the values of $\sum_{h \in [\ell]\setminus\{i,j\}}\alpha_h\ang{v_h, v_i}$ and $\sum_{h \in [\ell]\setminus\{i,j\}}\alpha_h\ang{v_h, v_j}$ are already determined by the fixed outcomes of $\alpha_h$ for $h\in [\ell]\setminus\{i,j\}$.

In the event $\cE_1(j,t)$, we have
\[|\ang{X_0,v_j} - \lambda_j|\le 10t\sqrt{\log m}.\]
So, the random variable $\alpha_j + \alpha_i\cdot A$ must be contained in the interval $I_j$ of length $20t\sqrt{\log m}/\rho_0$ centered around $\lambda_j/\rho_0 - \sum_{h \in [\ell]\setminus\{i,j\}}\alpha_h\ang{v_h, v_j}$. On the other hand, if $i$ is bad, we have 
\[|\ang{X_0, v_i} - \lambda_i| \le 10\sqrt{\log m},\]
so the random variable $\alpha_i + \alpha_j\cdot A$ must be contained in the interval $I_i$ of length $20\sqrt{\log m}/\rho_0$ centered around $\lambda_i/\rho_0 - \sum_{h \in [\ell]\setminus\{i,j\}}\alpha_h\ang{v_h, v_i}$. 

Thus, if $\cE_1(j,t)$ holds and $i$ is bad, taking the linear combination $(\alpha_j + A\cdot \alpha_i) - A(\alpha_i + A\cdot \alpha_j) = (1- A^2 )\alpha_j$, we see that $(1- A^2 )\alpha_j$ must be contained in the interval $I_j-AI_i$ of length at most $40t\sqrt{\log m}/\rho_0$ and so $\alpha_j$ must be contained in a particular interval of length at most $40t\sqrt{\log m}/((1-A^2)\rho_0)\le 400t\sqrt{\log m}/\rho_0$ (namely, the interval $(1- A^2 )^{-1}(I_j-AI_i)$). Here we used the inequality $1-A^2\ge 1/10$, which follows from $|A|\le 0.9$.

The probability that $\alpha_j\sim U[-1,1]$ lies in this interval is at most $200t\sqrt{\log m}/\rho_0$. Now conditioning on any such outcome of $\alpha_j$, for the event of $i$ being bad to occur, $\alpha_i\sim U[-1,1]$ must fall into an interval of length $20\sqrt{\log m}/\rho_0$ (which happens with probability at most $10\sqrt{\log m}/\rho_0$). Thus we can conclude that 
\[\P[i\text{ bad}\text{ and }\cE_1(j,t)\mid (\alpha_h)_{h\in [\ell]\setminus\{i,j\}}]\le  \frac{200t}{\rho_0}\sqrt{\log m}\cdot \frac{10}{\rho_0}\sqrt{\log m}=\frac{2000t}{\rho_0^2}\log m,\]
showing \eqref{eq:conditional-on-alpha-h}.

In the case where $S_j\le 400$, the desired inequality in the statement of the lemma follows directly from \eqref{eq:conditional-on-alpha-h}. Indeed, in this case we automatically have
\[\P[i\text{ bad}\text{ and }\cE_1(j,t)]\le  \frac{2000t}{\rho_0^2}\log m\le \frac{40000t}{\rho_0^2\sqrt{S_j}}\log m.\]

So let us assume $S_j > 400$ from  now on. We distinguish between the cases $\rho_0\le 10t\sqrt{\log m}$ and $\rho_0 > 10t\sqrt{\log m}$. In either case, in the event $\cE_1(j,t)$, we must have 
\[|\ang{X_0, v_j} - \lambda_j| = \Big|\rho_0\Big(\alpha_i+\alpha_j\cdot A + \sum_{h\in [\ell]\setminus\{i,j\}}\alpha_h\ang{v_h,v_i} \Big)- \lambda_j\Big| \le 10t\sqrt{\log m}.\]
So the sum $\rho_0\sum_{h\in [\ell]\setminus\{i,j\}}\alpha_h\ang{v_h,v_i}$ must lie in the interval of length $20t\sqrt{\log m}+4\rho_0$ centered at $\lambda_j$.

If $\rho_0\le 10t\sqrt{\log m}$, then this interval has length $20t\sqrt{\log m}+4\rho_0 \le 60t\sqrt{\log m}$. By \cref{lem:continuous-littlewood-offord}, the probability that $\rho_0\sum_{h\in [\ell]\setminus\{i,j\}}\alpha_h\ang{v_h,v_i}$ falls into this interval is at most 
\[\frac{60t\sqrt{\log m}}{\rho_0\sqrt{S_j - 1 - A^2}/\sqrt{3}}\le \frac{\sqrt{3}\cdot 60t\sqrt{\log m}}{\rho_0\sqrt{S_j - 2}}\le \frac{200t\sqrt{\log m}}{\rho_0\sqrt{S_j}}.\]
Conditioning on any such outcomes of $\alpha_h$ for all $h\in [\ell]\setminus\{i,j\}$, as well as on an outcome of $\alpha_j$, the event of $i$ being bad occurs only when $\alpha_i$ falls into a particular interval of length $20\sqrt{\log m}/\rho_0$, and the probability for this is at most $10\sqrt{\log m}/\rho_0$. So we can conclude that 
\[\P[i \text{ bad and }\cE_1(j,t)]\le \frac{200t\sqrt{\log m}}{\rho_0\sqrt{S_j}}\cdot \frac{10\sqrt{\log m}}{\rho_0} = \frac{2000t\log m}{\rho_0^2\sqrt{S_j}}.\]

If $\rho_0 > 10t\sqrt{\log m}$, in the event $\cE_1(j,t)$ the sum $\rho_0\sum_{h\in [\ell]\setminus\{i,j\}}\alpha_h\ang{v_h,v_i}$ must fall into the interval of length $20t\sqrt{\log m}+4\rho_0\le 6\rho_0$ centered at $\lambda_j$. By \cref{lem:continuous-littlewood-offord}, this happens with probability at most
\[\frac{6\rho_0}{\rho_0\sqrt{S_j - 1 - A^2}/\sqrt{3}}\le \frac{\sqrt{3}\cdot 6}{\sqrt{S_j - 2}}\le \frac{20}{\sqrt{S_j}}.\]
Conditioning on any such outcomes of $\alpha_h$ for all $h\in [\ell]\setminus\{i,j\}$, by \eqref{eq:conditional-on-alpha-h} the probability that $i$ is bad and $\cE_1(j,t)$ holds is at most  $(2000t/\rho_0^2)\log m$. So we can conclude 
\[\P[i\text{ bad}\text{ and }\cE_1(j,t)]\le  \frac{20}{\sqrt{S_j}}\cdot\frac{2000t}{\rho_0^2}\log m=\frac{40000t}{\rho_0^2\sqrt{S_j}}\log m. \qedhere\]
\end{proof}

\end{document}
